\section{Introduction}
\label{sec:intro}

In modern deep learning, the paradigm of training specialized, task-specific expert models and subsequently merging them into a single, unified multi-task model has gained significant traction \cite{ilharco2022editing, wortsman2022robust}. This approach, commonly known as parameter-space model merging, bypassed the prohibitive computational expense of joint multi-task training by performing algebraic or geometric alignment in the weight space. Prominent techniques such as Task Vectors \cite{ilharco2022editing} and TIES-Merging \cite{yadav2023ties} have demonstrated that parameter-space combinations can successfully combine diverse specialized capabilities with minimal performance degradation.

However, static weight-space merging methods are inherently limited. Because they apply a single, fixed set of merging coefficients across all test inputs, they cannot dynamically adapt the weight mixture to individual sample characteristics at inference time. To address this limitation, recent research has explored \emph{test-time dynamic model merging}, where a lightweight routing network maps the latent representations of input samples to layer-wise merging coefficients on the fly \cite{muqeeth2023l3, qwsmerge2025}. In theory, such dynamic routers allow a model to morph its parameters at test-time, activating relevant expert pathways based on the functional demands of the incoming sample.

However, as we empirically deconstruct in this paper, this theoretical promise is subject to a profound \emph{Dynamic Routing Paradox}: to avoid severe Vectorization Collapse, the dynamic router must be so heavily regularized (such as through zero-initialization and weight decay) that its learned coefficients are constrained to stay in a tight, high-entropy neighborhood of the static uniform compromise (achieving a Mean Absolute Deviation of only $0.0236$, or $2.36\%$, from the uniform $0.25$ baseline). This heavy regularization leaves the router with marginal functional flexibility, yielding a tiny $+1.16\%$ accuracy gain over naive, training-free Uniform Merging. This marginal benefit stands in stark contrast to the massive $O(B \cdot M)$ memory footprint expansion and reduced GPU arithmetic intensity of real-time parameter assembly, raising critical questions about the practical deployment viability of dynamic routing. We show that this paradox is not a unique artifact of linear merging; as we theoretically demonstrate in Appendix~\ref{app:nonlinear}, even if advanced non-linear alignment methods such as ZipIt or RegMean are deployed dynamically, they remain bound to the same fundamental paradox due to covariance estimation singular-value collapse in data-scarce splits.

Despite their theoretical appeal, current state-of-the-art dynamic model merging systems suffer from two critical, yet under-explored, vulnerabilities:
\begin{enumerate}
    \item \textbf{Sensitivity and Variance of Quantum-Inspired Architectures}: Recent work, such as Quantum Waveface Superposition (QWS-Merge) \cite{qwsmerge2025}, attempts to resolve routing complexities by adopting quantum-inspired wave-interference equations. However, as our empirically-driven audit reveals, these non-monotonic activation functions (e.g., wave cosine phase-interferometry) can introduce rugged, non-convex optimization landscapes. Under data-scarce 64-sample calibration splits, while QWS-Merge remains highly competitive under large batch sizes ($59.41\%$), it exhibits a higher variance across seeds and suffers a performance drop to $56.19\%$ (below the static Uniform baseline of $58.00\%$) under vectorized sample-wise deployment ($B=1$), highlighting the limits of complex activation dynamics in data-scarce regimes.
    \item \textbf{The Heterogeneity Collapse Confounder}: During deployment, inference streams are rarely homogeneous. When a dynamic router processes a heterogeneous batch containing mixtures of different tasks, the standard practice of batch-averaging dynamic layer-wise coefficients collapses the individual task coordinates back toward a uniform compromise. This batch-averaged "compromise weight" destroys the specialized adaptation of the model, a phenomenon we formalize as \emph{heterogeneity collapse}. For example, unregularized layer-wise classical linear routers experience a performance drop from 58.72\% joint accuracy in homogeneous streams to a lower 53.76\% under vectorized sample-wise deployment.
\end{enumerate}

To isolate these core routing and merging mechanics from confounding visual pre-training variables, we evaluate our methods entirely on a controlled, 192-dimensional synthetic Analytical Coordinate Sandbox. This allows us to systematically analyze routing behaviors in data-scarce (64 calibration samples) splits under true, customizable parameter-space overlap.

\begin{figure*}[t]
    \centering
    \begin{subfigure}[b]{0.48\textwidth}
        \centering
        \includegraphics[width=\textwidth]{fig1_lambda_sensitivity.png}
        \caption{Sensitivity curve over $\lambda_{var}$.}
        \label{fig:sensitivity}
    \end{subfigure}
    \hfill
    \begin{subfigure}[b]{0.48\textwidth}
        \centering
        \includegraphics[width=\textwidth]{fig2_heterogeneity_collapse.png}
        \caption{Deployment batch-size stress test.}
        \label{fig:heterogeneity}
    \end{subfigure}
    \caption{Empirical evaluation of our Prior-Driven Classical Routing Framework (specifically, zero-initialized Softmax routing with weight decay) and its variance-regularized variant (VR-Router). (a) Mapping the regularization sensitivity frontier of $\lambda_{var} \in [0.0, 10.0]$ across 10 independent random seeds. (b) Robustness under mixed-task deployment batch sizes $B \in \{1, 8, 32, 128, 512\}$. While unregularized standard routers collapse under vectorized sample-wise streaming, proper prior-layer designs (like VR-Router and L3\_Softmax\_WellReg) maintain highly stable, robust performance across all batch sizes.}
    \label{fig:main_results}
\end{figure*}

In this paper, we adopt a rigorously empirical and statistical methodology—which we call \emph{The Empiricist} approach—to dissect these challenges and propose a mathematically simpler, yet profoundly more robust solution. Our key finding is that the stability of dynamic routing in data-scarce splits depends fundamentally on proper initialization and weight decay, rather than on complex activation functions or explicit loss penalties alone.

Our proposed framework, which we call \textbf{Prior-Driven Classical Routing} (or \textbf{Zero-Initialized Softmax Routing}), utilizes standard, zero-initialized Softmax routing layers regularized with $L_2$ weight decay. We project high-dimensional latent representations onto a low-dimensional unit sphere via normalized random projections to prevent overfitting. We then map these representations to layer-wise coefficients via classical linear projections. To combat heterogeneity collapse, we perform true sample-specific vectorized assembly during calibration. Within this framework, we also explore an optional group-level \textbf{Task-Variance Regularization ($\mathcal{L}_{VR}$)} penalty (which we formulate as VR-Router) that minimizes the intra-task variance of predicted coefficients. Through a systematic baseline audit, we show that standard L3-Softmax---when properly zero-initialized and regularized with weight decay---performs identically to VR-Router, proving that simple architectural priors are the true, necessary drivers of generalization.

To validate our claims, we do not rely on selective seeds or isolated hyperparameter settings. Instead, we execute a massive empirical evaluation pipeline:
\begin{itemize}
    \item We perform a \textbf{Statistical Significance Audit} across 10 independent random seeds, completely reconstructing the visual representation sandbox with distinct feature spaces and calibration splits for each.
    \item We map the \textbf{Regularization Sensitivity Frontier} by sweeping the variance penalty weight $\lambda_{var} \in [0.0, 10.0]$ across 10 values, pinpointing the optimal boundary.
    \item We conduct an \textbf{Inference Stream Heterogeneity Stress Test}, sweeping deployment batch sizes $B \in \{1, 8, 32, 128, 512\}$ to verify batch-independent robustness.
    \item We present an \textbf{Exhaustive Ablation Study} of all loss components ($\mathcal{L}_{CE}, \mathcal{L}_{reg}, \mathcal{L}_{VR}$) to empirically isolate the exact drivers of generalization.
\end{itemize}

Our results are definitive and reveal a crucial, previously unstudied phenomenon: the \emph{Batch-Average Smoothing Confounder}. Under standard evaluation on heterogeneous streams with batch size $B=256$, standard dynamic routers (such as random-initialized L3-Softmax) appear to perform well (59.35\% $\pm$ 1.33\%), but this is a false artifact. The batch-averaging of predicted coefficients acts as an implicit smoothing operator that masks the severe overfitting of the router on the 64-sample calibration set. When deployed in true, real-time batch-independent or sample-wise vectorized pipelines ($B=1$), this smoothing mask is removed, causing a severe \emph{Vectorization Collapse} where random-initialized L3-Softmax drops to \textbf{41.09\% $\pm$ 3.73\%} (nearly 17\% below Uniform Merging). In contrast, our well-regularized, zero-initialized L3-Softmax baseline (\texttt{L3\_Softmax\_WellReg}) and VR-Router completely resolve this vulnerability. They achieve stable, highly-generalizing joint accuracies of \textbf{59.16\% $\pm$ 1.17\%} and \textbf{59.14\% $\pm$ 1.18\%} respectively under both batch-average ($B=256$) and vectorized sample-specific ($B=1$) deployments, outperforming Uniform Merging (58.00\% $\pm$ 1.13\%) and all unregularized dynamic baselines. This proves that proper routing-layer regularization and initialization are the true, foundational drivers of deployment stability.
